\scanpage{24}
% L01 heading
Invariant\unsure{s} \deleted{\unsure{des}}d'un automorphismes d\deleted{e}'une droite projective.

% L02–06
a) Soit $G$ une forme de $\mathrm{GP}(1)$, ou de $\mathrm{Sl}(2)$ \inserted{sur $S$}, -- i.e. un schéma en groupes $\frac12$ simple de rang $1$ sur schéma $S$. Considérons le foncteur
\[
\mathbb I(G)=\underline{\mathrm{Hom}}_{\rm sch}(G,\mathbb E^1_S)^{\unsure{\leftrightarrows|\leftrightarrow|\circlearrowleft}}
\overset{\text{déf}}{=}\underline{\mathrm{Hom}}_{\rm sch}(G,\mathbb E^1_S)^G
\]
où $G$ opère sur le $\underline{\mathrm{Hom}}$ \deleted{via} autom. int. sur $G$.

% L07–14
Ce foncteur est représentable, \inserted{(vrai pour tt $G$ $\frac12$ simple)} et on a un isomorphisme canonique
\[
\mathbb I(G)\xrightarrow{A}\mathbb E^1_S
\]
qui, en tant que morphisme, peut être caractérisé par la propriété suivante : Si $G_m\overset{i}{\lhook\joinrel\longrightarrow}G$ est un monomorphisme, alors
\[
A(i(\lambda))=\lambda+\lambda^{-1}
\]
et itou après tout chgt de base.

% L15–21
\deleted{\uline{Proposition} Soient $u,v$ deux sections de $G$, telles que \inserted{on suppose} \inserted{\unsure{a) Donc que $u$|2) Prouvons que}}}
\deleted{\unsure{ou}} On a donc
\[
A(1)=2.
\]
[ \deleted{Supposons $G$ forme de $\mathrm{GP}(1)$, et} \inserted{Soit} \deleted{et si} $\sigma$ \deleted{est} une section de $G$ sur $S$ telle que $\sigma^2=\mathrm{id}$, $\sigma_s\ne\mathrm{id}_s\quad\forall s\in S$, \deleted{$\sigma$ soit d'ordre $2$ \unsure{exact} (donc ce qu'on peut appeler une réflexion de $G$)} alors
\[
A(\sigma)=-2
\]
% slanted
\margin{\uline{NB} Dans le cas d'une forme de $\mathrm{GP}(1)$, $\sigma$ est une ``réflexion'', dans le cas de $\mathrm{Sl}(2)$, en car rés. $\ne2$, \unsure{c'est la} section génératrice du centre\ldots}

% L22–26
C'est clair si $2$ inv. sur $S$, car alors $\sigma$ est contenu dans un \deleted{(unique)} tore maximal, \inserted{(unique dans le cas de $\mathrm{GP}(1)$ -- dans le cas de $\mathrm{Sl}(2)$, $\sigma$ est central)} qu'on peut supposer scindé (par localisation), et on a alors $\lambda=-1$, d'où $A(u)=-1+1/{-1}=-2$. Pour le cas général, on fait un argument de coeff universels \ldots. ]

% L27–28
\deleted{\uline{Prop} Dire que} Supposons que $u$ soit \uline{unipotent}, par quoi on entend qu'il existe loc. fpqc une \deleted{forme}
% archive sheet number: 23
